% !TEX root =  Lecture18.tex


%%%%%%%%%%%%%%%%%%%%%%%%%%%%%%%%%%%%
% Recap/Agenda
%%%%%%%%%%%%%%%%%%%%%%%%%%%%%%%%%%%%
\begin{frame}
\frametitle{Module 4: Bayesian Statistics}
    \begin{itemize}
        \item \hl{Previously:} Lecture 17 opened this module: priors, Bayes' rule, sequential updating.
        \item \hl{This time:} Lecture 17 spread a prior over five candidate values of $p$. Now the grid becomes a \hl{curve}.
          \begin{itemize}
              \item the \hl{Beta} prior, and why Beta $+$ Binomial hands back a Beta;
              \item what makes a prior \hl{conjugate}, and the one pattern all three models share;
              \item two more families: \hl{Normal-Normal} and \hl{Gamma-Poisson};
              \item how much your prior is worth, measured in observations, and where conjugacy runs out.
          \end{itemize}
        \item \hl{Reading:} Bayes Rules!\ Chapters 3 and 5.
        \item \hl{Homework:}
    \end{itemize}
\end{frame}


%%%%%%%%%%%%%%%%%%%%%%%%%%%%%%%%%%%%
% Learning objectives:
%%%%%%%%%%%%%%%%%%%%%%%%%%%%%%%%%%%%
\begin{frame}
    \frametitle{Lecture Objectives}
    \goals{%
    \begin{itemize}
    \item explain what makes a prior \hl{conjugate} to a likelihood, and derive the Beta-Binomial posterior by multiplying shapes, with no integration;
    \item apply all three update rules (\hl{Beta-Binomial}, \hl{Normal-Normal}, \hl{Gamma-Poisson}) and read their parameters as ``prior successes'' or ``prior observations'';
    \item state the \hl{prior sample size} $n_0$ for each family, and use it to compute the weights on prior and data;
    \item explain why the posterior mean is always a \hl{weighted average} of the prior mean and the MLE, and predict which side it falls on;
    \item say what conjugate models buy you, and \hl{where they run out}.
    \end{itemize}}
\end{frame}


%%%%%%%%%%%%%%%%%%%%%%%%%%%%%%%%%%%%
\begin{frame}
    \frametitle{Lecture Objectives}
    \objectives{%
    \begin{itemize}
    \item \textbf{(M4, LO3)} Conjugate Models
    \end{itemize}}
    \vspace{0.3em}
    {\small Building on, and revisiting:}
    \objectives{%
    \begin{itemize}
    \item \textbf{(M4, LO1)} Bayesian Inference: the prior, likelihood and posterior of Lecture 17, now with the parameter on a continuous scale
    \item \textbf{(M4, LO2)} Comparing Bayesian and Frequentist Approaches: the MLE reappears here as one end of a tug-of-war
    \end{itemize}}
\end{frame}
